%=========================================================
%               CLASSE DU DOCUMENT
%=========================================================
\documentclass[12pt,a4paper]{article}

%=========================================================
%         ENCODAGE ET POLICES DE CARACTÈRES
%=========================================================
\usepackage[utf8]{inputenc}    % Encodage UTF-8 (inutile avec XeLaTeX/LuaLaTeX)
\usepackage[T1]{fontenc}       % Meilleure gestion des caractères accentués
\usepackage[french]{babel} % <--- Ce package fournit \og et \fg{}
\usepackage{lmodern}           % Police Latin Modern
\usepackage{times}             % Police Times (à supprimer si tu préfères Latin Modern)

%=========================================================
%                MATHÉMATIQUES
%=========================================================
\usepackage{amsmath}           % Environnements mathématiques avancés
\usepackage{amssymb}           % Symboles mathématiques supplémentaires
\usepackage{amsfonts}          % Polices mathématiques
\usepackage{mathrsfs}          % Police calligraphique (\mathscr)
\usepackage{tkz-tab}           % Tableaux de signes et de variations
\everymath{\displaystyle}      % Affiche toutes les formules en style "display"

%=========================================================
%               MISE EN PAGE
%=========================================================
\usepackage[
    left=1cm,
    right=1cm,
    top=1cm,
    bottom=1.7cm
]{geometry}                    % Marges du document

\usepackage{multicol}          % Texte sur plusieurs colonnes
\usepackage{float}             % Placement précis des figures et tableaux

%=========================================================
%            TABLEAUX ET LISTES
%=========================================================
\usepackage{array}             % Colonnes de tableaux personnalisées
\usepackage{multirow}          % Fusion de lignes dans les tableaux
\usepackage{enumitem}          % Personnalisation des listes

\usepackage{tikz}
\usetikzlibrary{angles,quotes,arrows.meta}
%=========================================================
%            COULEURS ET BOÎTES
%=========================================================
\usepackage[most]{tcolorbox}   % Boîtes colorées très personnalisables
\tcbuselibrary{skins,hooks}    % Bibliothèques supplémentaires de tcolorbox

\usepackage{mdframed}          % Cadres autour de paragraphes
\usepackage{varwidth}          % Largeur automatique des boîtes

%=========================================================
%                DESSINS ET GRAPHIQUES
%=========================================================
\usepackage{tikz}              % Dessins vectoriels
\usetikzlibrary{calc}
\usetikzlibrary{
    arrows,
    shapes.geometric,
    fit,
    patterns
}

\usepackage{pgfplots}          % Tracés de fonctions et graphiques
\pgfplotsset{compat=1.15}

%=========================================================
%              LIENS HYPERTEXTES
%=========================================================
\usepackage[
    colorlinks=true,
    linkcolor=blue,
    urlcolor=blue,
    citecolor=blue
]{hyperref}

%=========================================================
%          ÉLÉMENTS EN ARRIÈRE-PLAN
%=========================================================
\usepackage{eso-pic}           % Ajout de filigranes ou d'éléments sur chaque page

%=========================================================
%              SYMBOLES DIVERS
%=========================================================
\usepackage{pifont}            % Symboles (✓, ✗, etc.)

\DeclareUnicodeCharacter{2460}{\text{\ding{172}}} % Pour ①
\DeclareUnicodeCharacter{2461}{\text{\ding{173}}} % Pour ②
%=========================================================
%      PERSONNALISATION DES LISTES NUMÉROTÉES
%=========================================================

\newcommand\circitem[1]{%
\tikz[baseline=(char.base)]{
\node[circle,draw=gray,fill=red!55,
minimum size=1.2em,inner sep=0] (char) {#1};}}

\newcommand\boxitem[1]{%
\tikz[baseline=(char.base)]{
\node[fill=cyan,
minimum size=1.2em,inner sep=0] (char) {#1};}}

\setlist[enumerate,1]{label=\protect\circitem{\arabic*}}
\setlist[enumerate,2]{label=\protect\boxitem{\alph*}}

%=========================================================
%            COMMANDES PERSONNALISÉES
%=========================================================

% Soulignement coloré
\newcommand{\myul}[2][black]{%
\setulcolor{#1}\ul{#2}\setulcolor{black}}

% Tabulation horizontale
\newcommand\tab[1][1cm]{\hspace*{#1}}

%=========================================================
%             COMPTEURS PERSONNALISÉS
%=========================================================

%---------- Exemple ----------
\newcounter{exemple}

\newcommand{\exemple}{%
\refstepcounter{exemple}
\textbf{\textcolor{orange}{Exemple \theexemple : }}
\ignorespaces}

%---------- Solution ----------
\newcounter{solution}

\newcommand{\solution}{%
\refstepcounter{solution}
\textbf{\textcolor{orange}{Solution \thesolution : }}
\ignorespaces}

%---------- Exercices ----------
\definecolor{myorange}{rgb}{1.0,0.8,0.0}

\newcounter{exerciceapp}

\newcommand{\exerciceapp}{%
\refstepcounter{exerciceapp}
\textbf{\textcolor{myorange}
{Exercice d'application \theexerciceapp : }}
\ignorespaces}

%---------- Corrections ----------
\newcounter{correction}

\newcommand{\correction}{%
\refstepcounter{correction}
\textbf{\textcolor{myorange}
{Correction \thecorrection : }}
\ignorespaces}

%---------- Remarques ----------
\definecolor{myorange1}{rgb}{1.0,1.0,0.0}
\usepackage{enumitem}

\definecolor{myred}{RGB}{180, 40, 40}
\definecolor{myblue}{RGB}{20, 50, 120}
\newcounter{remarque}

\newcommand{\remarque}{%
\refstepcounter{remarque}
\textbf{\textcolor{myorange1}
{Remarque \theremarque : }}
\ignorespaces}

%=========================================================
%                 FILIGRANE
%=========================================================

\AddToShipoutPicture{
    \AtTextCenter{
        \makebox[0pt]{
            \rotatebox{45}{
                \textcolor[gray]{0.6}{
                    \fontsize{5cm}{5cm}\selectfont PGB
                }
            }
        }
    }
}

\begin{document}
% En-tête personnalisée
\begin{center}
    \Large\textbf{\underline{\textcolor{red}{Polynômes}}}\\[-0.1cm]
    \normalsize\textbf{Prof : M. BA} \hfill \textbf{Classe : 1erS2}\\[-0.1cm]
    \textbf{Année scolaire : 2025 -- 2026}
\end{center}

\section*{Objectifs spécifiques}
\begin{itemize}[label=$\checkmark$]
    \item Diviser un polynôme \( P(x) \) par \( x - a \) lorsque \( P(a) = 0 \);
    \item Diviser un polynôme \( P(x) \) par \( x - a \), puis le quotient \( Q(x) \) par \( x - b \) lorsque \( P(a) = P(b) = 0 \);
    \item Factoriser un polynôme;
    \item Étudier le signe d'un polynôme;
    \item Étudier le signe d'une fraction rationnelle.
\end{itemize}

\section*{Prérequis}
\begin{itemize}[label=$\checkmark$]
    \item Factorisation de polynômes de degré \( \leq 4 \) par identification et division euclidienne;
    \item Résolution d'inéquations de degré \( \leq 4 \).
\end{itemize}

\section*{Supports didactiques}
\begin{itemize}[label=$\checkmark$]
    \item Collection Hachettes classiques \( 1^{\text{ère}} A_1 \) et \( B \);
    \item CIAM littéraire \( 1^{\text{ère}} \);
    \item Ordinateur.
\end{itemize}

\section*{Plan du chapitre}
\begin{enumerate}
    \item \textbf{Rappels}
    \begin{enumerate}
        \item Factorisation d'un trinôme du second degré
        \item Signe d'un trinôme du second degré
    \end{enumerate}
    \item \textbf{Factorisation d'un polynôme de degré \( n \geq 3 \) connaissant une racine}
    \begin{enumerate}
        \item Par la division euclidienne
        \item Par la méthode d'identification
        \item Par la méthode de Horner
        \begin{enumerate}
            \item Exemple
            \item Exercice d'application
        \end{enumerate}
    \end{enumerate}
    \item \textbf{Signe d'un polynôme et résolution d'inéquations}
    \begin{enumerate}
        \item Exemple
        \item Exercice d'application
    \end{enumerate}
    \item \textbf{Fraction rationnelle}
    \begin{enumerate}
        \item Définition et exemples
        \item Signe d'une fraction rationnelle
    \end{enumerate}
\end{enumerate}

\newpage

\section*{I. Rappels}

\subsection*{1. Factorisation d'un trinôme du second degré}

Le tableau suivant permet de factoriser un trinôme du second degré \(ax^2 + bx + c\).

\begin{table}[h]
\centering
\begin{tabular}{|c|c|}
\hline
\textbf{Signe du discriminant} \(\Delta = b^2 - 4ac\) & \textbf{Factorisation du trinôme} \(ax^2 + bx + c\) \\
\hline
\(\Delta < 0\) & \(ax^2 + bx + c\) ne peut pas se factoriser. \\
\hline
\(\Delta = 0\) & \(ax^2 + bx + c = a(x - x_0)^2\) où \(x_0 = -\dfrac{b}{2a}\) \\
\hline
\(\Delta > 0\) & \(ax^2 + bx + c = a(x - x_1)(x - x_2)\) \\
 & où \(x_1 = \dfrac{-b + \sqrt{\Delta}}{2a}\) et \(x_2 = \dfrac{-b - \sqrt{\Delta}}{2a}\) \\
\hline
\end{tabular}
\end{table}

\paragraph{Exemples :} Factorisons les trinômes du second degré suivants :
\[
f(x) = x^2 - 2x + 3 \quad;\quad g(x) = -9x^2 + 6x - 1 \quad;\quad h(x) = x^2 - x - 6
\]

\subsection*{2. Signe d'un trinôme du second degré}

Le signe d'un trinôme du second degré s'obtient généralement à l'aide d'un tableau de signe.

\subsubsection*{Si \(\Delta < 0\)}
Alors le trinôme n'a pas de racine et son tableau de signe est le suivant :

\begin{center}
    \begin{tikzpicture}[node style/.style={fill opacity=0,text opacity=1}]
    \tkzTabInit[
        lgt=3.5,
        espcl=2.5
    ]{$x$/.5,$ax^2 + bx + c$/.7}
    {$-\infty$,$+\infty$}
    \tkzTabLine{,\text{signe(a)},}
    \end{tikzpicture}
\end{center}

\paragraph{Exemple :} Étudions le signe de \(-x^2 + x - 2\).

\subsubsection*{Si \(\Delta = 0\)}
Alors le trinôme a une seule racine dite racine double $x_0$ et son tableau de signe est le suivant :

\begin{center}
    \begin{tikzpicture}[node style/.style={fill opacity=0,text opacity=1}]
    \tkzTabInit[
        lgt=3.5,
        espcl=2.5
    ]{$x$/.5,$ax^2 + bx + c$/.7}
    {$-\infty$,$x_0$,$+\infty$}
    \tkzTabLine{,\text{signe(a)},z,\text{signe(a)},}
    \end{tikzpicture}
\end{center}

\paragraph{Exemple :} Étudions le signe de \(9x^2 - 6x + 1\).

\subsubsection*{Si \(\Delta > 0\)}
Alors le trinôme a deux racines distinctes \(x_1\) et \(x_2\) (avec \(x_1 < x_2\)) et son tableau de signe est le suivant :

\begin{center}
    \begin{tikzpicture}[node style/.style={fill opacity=0,text opacity=1}]
    \tkzTabInit[
        lgt=3.5,
        espcl=2.5
    ]{$x$/.5,$ax^2 + bx + c$/.7}
    {$-\infty$,$x_1$,$x_2$,$+\infty$}
    \tkzTabLine{,\text{signe(a)},z,-\text{signe(a)},z,\text{signe(a)},}
    \end{tikzpicture}
\end{center}

\section*{II. Factorisation d'un polynôme de degré \( n \geq 3 \) connaissant une racine}

\subsection*{a. Par la division euclidienne}

\paragraph{Exemple :} Soit le polynôme \( x^3 - 2x^2 - 5x + 6 \)
\begin{enumerate}
    \item Vérifier que \(-2\) est une racine de ce polynôme.
    \item Factoriser ce polynôme en utilisant la méthode de la division euclidienne.
\end{enumerate}

\subsection*{b. Par la méthode d'identification}

\paragraph{Rappels :}
\begin{itemize}
    \item Si \( P(x) \) est un polynôme de degré 2 alors \( P(x) \) peut s'écrire sous la forme 
    \[
    P(x) = ax^2 + bx + c
    \]
    avec \( a, b, c \) des réels constants.
    \item Si \( P(x) \) est de degré 3 alors \( P(x) \) peut s'écrire sous la forme 
    \[
    P(x) = ax^3 + bx^2 + cx + d
    \]
    avec \( a, b, c, d \) des réels constants.
\end{itemize}

\paragraph{Exemple :}
\begin{enumerate}
    \item Vérifier que 1 est racine de \( x^3 - 7x + 6 \).
    \item Factoriser ce polynôme en utilisant la méthode d'identification.
\end{enumerate}

\subsection*{c. Par la méthode de Horner}

\paragraph{Exemple :} Soit le polynôme \( x^3 - 2x^2 - 5x + 6 \)
\begin{enumerate}
    \item Vérifier que \(-2\) est racine de \( x^3 - 2x^2 - 5x + 6 \).
    \item Factoriser \( x^3 - 2x^2 - 5x + 6 \) en utilisant la méthode de Horner.
\end{enumerate}

\[
\begin{array}{c|cccc}
\text{Coefficients de } P(x) \text{ dans l'ordre décroissant des puissances} & 1 & -2 & -5 & 6 \\
\hline
\text{Racine : } (-2) & & & & \\
\text{Coefficients de } Q(x) \text{ dans l'ordre décroissant des puissances} & & & & \\
\hline
 & 1 & -4 & 3 & 0
\end{array}
\]

\[
x^3 - 2x^2 - 5x + 6 = (x + 2)(x^2 - 4x + 3) + 0
\]

On remarque que \( x^2 - 4x + 3 = (x - 1)(x - 3) \). Par suite :
\[
x^3 - 2x^2 - 5x + 6 = (x + 2)(x - 1)(x - 3)
\]

\paragraph{Exercice d'application :}
\begin{enumerate}
    \item Vérifier que 1 est racine de \( x^3 - 7x + 6 \).
    \item Factoriser ce polynôme en utilisant la méthode de Horner.
\end{enumerate}
\end{document}